\documentclass[11pt]{article}
\usepackage{amsmath}
\usepackage{graphicx}
\usepackage{booktabs}
\usepackage{array}
\newcolumntype{L}[1]{>{\raggedright\arraybackslash}p{#1}}
\usepackage[numbers,sort&compress]{natbib}
\usepackage[hidelinks]{hyperref}
\def\UrlBreaks{\do\/\do-\do_\do.}
\Urlmuskip=0mu plus 1mu
\emergencystretch=3em

\newcommand{\pending}[1]{\textbf{Results: pending} (produced by \nolinkurl{#1}).}
\newcommand{\internal}[1]{internal planning note \nolinkurl{#1}, unpublished}

\title{Governing Public World Simulators:\\Misuse, Consent, Labeling and Uncertainty in a Forkable Simulation of Society}
\author{Real World Research\\\small Author names and affiliations to be added before submission\\\small Position paper, revised draft of 11 October 2026 (internal review round 2)}
\date{}

\begin{document}
\maketitle

\begin{abstract}
A public world simulator is a planned system that would hold a synthetic population, built to resemble real populations, as numeric records grouped into segments; use a language model to read the daily news into bounded changes in those segments' beliefs; publish forecasts about real events; and let anyone fork the simulated world, inject an event, and read what the simulated population ``does''. It differs from a private synthetic-respondent panel in three ways that matter for governance: its outputs are public and can borrow the authority of a scored forecaster, anyone can query it, and it speaks about real societies, real elections and, if permitted, real people. We argue that such a system should be governed as a \emph{forecasting instrument under public scoring}, not as a substitute for human respondents and never as a tool for testing or targeting political persuasion. For the persuasion harm specifically, we argue (but do not show) that structural limits in the output schema and API work better than content filtering alone. We give a threat model for five harms: persuasion optimization, simulating named people, election and political-publication harms, false authority from flattened synthetic populations, and data-rights and group-privacy failures. We propose rules for each, for one planned instance (Real World); the company has not yet adopted them, and we say where they conflict with its commercial plans. Because none of the rules has been tested, we specify seven audit protocols, to be pre-registered by hash before any data are collected, covering leakage of persuasion rankings against a self-built persona-panel baseline, the event-text classifier, label persistence, display comprehension and the transfer of authority to unscored outputs, re-identification and attribute disclosure, reliance on validation grades, and conformance checks including the market-blind scoring run. Each protocol states its unit, sample size, metric, test, decision rule with a candidate value for every free parameter, its power against a system that is in fact safe, its retest limit, and the changes that invalidate a pass. All results are pending.
\end{abstract}

\section{Introduction}

Language models conditioned on demographic backstories can reproduce some group-level survey distributions \citep{argyle2023outofone}, and agents built from two-hour interviews of 1,052 Americans predicted held-out General Social Survey answers at 83\% of the participants' own two-week test-retest consistency, against 74\% for agents given demographics alone \citep{park2024selfreports}. Strong rebuttals followed: synthetic respondents compress variance, give different answers to the same prompt over a three-month period, and get regression coefficients wrong \citep{bisbee2024synthetic}; individual ``digital twins'' correlate weakly with the person they model (average $r=0.20$ across 164 outcomes in 19 pre-registered studies) \citep{peng2025funhouse}; and default model opinions sit closest to a small set of countries \citep{durmus2023global,tao2023cultural}. In a September 2026 test by Pew Research Center, the best configuration missed human survey benchmarks by 11.2 to 14.7 percentage points on average, depending on the wave \citep{pew2026silicon,pew2026report}.

Most of the ethics discussion so far concerns \emph{private} uses of simulated respondents: a firm replaces part of a survey, or a lab replaces human participants \citep{agnew2024illusion,wang2024misportray}. This paper concerns a different object, which we call a \emph{public world simulator}, and one planned instance of it, Real World. No product code exists at the time of writing; Section~\ref{sec:system} describes the design as planned, in the conditional. In brief, the planned engine would hold about 100,000 synthetic agents as numeric records grouped into about 1,000 to 2,000 segments in total; each agent's belief is its segment's position plus a fixed individual offset, and the model is called per news story, per question or per segment, never per agent. The agent count is therefore not a sample size, and it is not a count of language-model agents. Each morning a data-assimilation cycle would ingest news, polls and search trends, correct segment beliefs toward the polls, and treat news and search as bounded forcing; prediction-market prices are never assimilated and serve only as the scoring baseline. Forecasts would be published as hash commitments before their outcomes exist and scored in public with the Brier score, against the market price at the scheduled call time as a baseline. Anyone could fork the day's world, add an event (``a major chip export ban''), and run it forward for seven simulated days.

Three properties set a public simulator apart. It is \emph{public}: what it says can be screenshotted, quoted and repeated without the caveats around it. It is \emph{open to queries}: anyone can try an input, including inputs designed to extract something harmful. And it carries \emph{borrowed authority}: a scored public track record is the system's strongest claim to trust, and that trust can leak to outputs that were never scored, such as counterfactual forks.

\paragraph{Position.} We argue that a public world simulator should be governed as a forecasting instrument under public scoring. Concretely: (i) it simulates populations, never named persons; (ii) users act on the world through \emph{shocks} (changes in the state of the world that happen to everyone), never through \emph{messages} (communication aimed at a population whose response is then compared), and it is not offered for political persuasion testing in any form; (iii) every output is labeled as a simulation in the machine-readable schema as well as on the page; (iv) every number carries an uncertainty statement that says what it is, every forecast carries the system's track record with its sample size, and unscored outputs say that they are unscored; (v) checks that gate a capability fail closed; and (vi) each rule is paired with an audit or conformance check, pre-registered before it runs, whose result is published whether it flatters the system or not.

\paragraph{Contributions.}
\begin{enumerate}
  \item A threat model for public world simulators covering five harms and the affordances that enable each, with the marginal-risk question made explicit (Sections~\ref{sec:system}--\ref{sec:threats}).
  \item A set of concrete proposed rules for Real World, with the reasoning, the legal or platform rule each answers to, the residual risk each leaves, and the company decisions they depend on, none of which has been made (Section~\ref{sec:policies}, Table~\ref{tab:policies}).
  \item A consent framework for three groups whose consent is easy to ignore: people who might later opt into a personal digital twin, the survey respondents whose records ground the synthetic population, and the users who answer the system's daily questions (Section~\ref{sec:consent}).
  \item Seven audit protocols, to be pre-registered by hash and, for human-subjects studies, on a public registry, before any data are collected. Each has a unit, sample size, metric, test, decision rule with candidate parameter values, power against a safe system, a retest limit and re-audit triggers, and is mapped to the release stage it gates (Section~\ref{sec:protocols}, Appendix~\ref{app:protocols}). All results are pending.
  \item An account of where this position may be wrong (Section~\ref{sec:limits}).
\end{enumerate}

This paper leaves the technical content to companion papers, all drafts in preparation: how the population is built (\emph{Unequally Known}), how the filter that corrects beliefs toward observations can be upgraded (\emph{Belief Assimilation}), how dispersion is corrected (\emph{Correcting Variance Collapse}), how forecasts are committed and scored (\emph{Receipts}), how the system is benchmarked (\emph{The Real World Benchmark}), the system as a whole (\emph{Don't Ask a Model, Ask a World}), and what the silicon-sample literature supports (\emph{Silicon Samples, 2022--2026}). We refer to them only where a governance choice depends on a technical one. Where a statement rests on the company's internal planning notes, we cite the note by its path in the company's research repository and mark it as unpublished.

\section{Background and Related Work}

\paragraph{Silicon samples and their failure modes.} \citet{argyle2023outofone} introduced ``silicon samples'' and the idea of algorithmic fidelity. \citet{bisbee2024synthetic} found that group means were close to the American National Election Studies but with less variation, regression estimates often differed, and the same prompt gave significantly different results over a three-month period. \citet{santurkar2023whose} showed misalignment with many US demographic groups that demographic steering did not fix. \citet{durmus2023global} found default model responses closest to the opinions of the USA and some European and South American countries, and that prompting for a country's perspective can reproduce harmful stereotypes; \citet{tao2023cultural} found cultural values resembling English-speaking and Protestant European countries in five models, with cultural prompting improving alignment in 71--81\% of countries and territories for recent models. Alignment tuning has mixed effects on such gaps and can widen them \citep{ryan2024unintended}. A review found that fewer than half of LLM social-simulation studies test variance at all, and most of those that do find it lower than in human populations \citep{wu2025boundary}. \citet{anthis2025promising} argue the method is still promising if these challenges are addressed. These findings ground a governance claim, not only a methods claim: a system with known distortions has a duty not to present its outputs as more certain or more individual than they are.

\paragraph{Predicting experiments and events.} \citet{ashokkumar2026predict} prompted GPT-4 to simulate representative US samples in 70 pre-registered survey experiments and found predicted treatment effects strongly correlated with measured ones, at about the accuracy of pooled human forecasters, while systematically overestimating effect sizes. This is the closest evidence that a simulator's comparisons between treatments can carry real signal, and it bears directly on the persuasion harm below. For event forecasting, a retrieval-augmented language-model system approached the aggregate of competitive human forecasters \citep{halawi2024forecasting}, and ForecastBench provides a continuously updated benchmark against which such systems are compared \citep{karger2025forecastbench}. A same-feed language model, one reading the same news with no population, is therefore the natural baseline for any claim that a simulated population adds forecasting skill, and it is one of the baselines we commit to below.

\paragraph{Replacing participants and representational harm.} \citet{wang2024misportray} argue, and show with 3,200 human participants across 16 identities and four models, that LLMs misportray and flatten identity groups, and that mitigations reduce but do not remove the harm. \citet{agnew2024illusion} argue that substitution proposals conflict with the values of representation, inclusion and understanding that human-participant research exists to serve. A world simulator that speaks \emph{for} a country's population inherits both critiques at national scale.

\paragraph{Digital twins of individuals.} Interview-grounded agents \citep{park2024selfreports} and large twin datasets \citep{toubia2025twin2k} make individual simulation concrete. Ethicists have responded with frameworks for ``digital duplicates'' built on harm and proportionality, information permission, transparency, authenticity and appropriate access \citep{hurshman2026duplicates}; with an argument that the ``digital twin'' metaphor itself mischaracterizes personalized LLMs and fosters misplaced trust \citep{annoni2026metaphor}; and with a critique of digital cloning in research drawing on feminist theories of agency \citep{brooke2025clones}.

\paragraph{Persuasion, microtargeting and marginal risk.} In three experiments with 76,977 participants, 19 LLMs and 707 political issues, persuasion post-training and prompting raised persuasiveness by up to 51\% and 27\% respectively, more than personalization or scale did, and the same methods lowered factual accuracy \citep{hackenburg2025levers}. The evidence on tailoring is mixed. In a pre-registered debate experiment ($N=900$), GPT-4 with access to participants' sociodemographic data produced an 81.2\% relative increase in the odds of higher post-debate agreement compared with human opponents, in the debate pairs where AI and humans were not equally persuasive \citep{salvi2024persuasiveness}, and psychologically matched messages outperformed generic ones in \citet{matz2024personalized} and \citet{simchon2024microtargeting}. In contrast, a pre-registered experiment ($n=8{,}587$) found microtargeted GPT-4 messages no more persuasive in aggregate than non-targeted ones \citep{hackenburg2024microtargeting}, and \citet{hackenburg2025levers} found personalization a weak lever. \citet{goldstein2023influence} map the language-model-to-influence-operation pipeline and conclude that no reasonable mitigation can be expected to prevent the threat fully, while a combination of mitigations may make an important difference. The dual-use lesson of \citet{urbina2022dualuse} applies: a simulator that estimates how a population responds can be turned toward choosing what to say to it. Against all of this, \citet{kapoor2024open} propose that misuse claims be judged by \emph{marginal} risk over pre-existing technology; we apply that test in Section~\ref{sec:threats} and build it into the E1 baseline.

\paragraph{Privacy of synthetic data and statistical releases.} Synthetic data is not anonymous by construction: across many generative models it either fails to prevent inference attacks or loses utility, with a privacy gain that is hard to predict \citep{stadler2022groundhog}. Population uniqueness, not sample uniqueness, determines re-identification risk; a generative model estimated that 99.98\% of Americans would be correctly re-identified in any dataset using 15 demographic attributes \citep{rocher2019reidentification}. Language models can emit verbatim training data \citep{carlini2021extracting}. The disclosure-control literature supplies the vocabulary we use: $k$-anonymity requires each released record to be indistinguishable from at least $k-1$ others on its quasi-identifiers \citep{sweeney2002kanonymity}, and $\ell$-diversity addresses the attribute disclosure that $k$-anonymity leaves open when the records in a cell share a sensitive value \citep{machanavajjhala2007ldiversity}. Aggregate query interfaces leak small cells through overlapping queries, the ``tracker'' attack \citep{denning1979tracker}, and answering too many aggregate queries too accurately allows the underlying records to be reconstructed \citep{dinur2003revealing}. That result, demonstrated against published census tables, is why the US Census Bureau moved its 2020 disclosure avoidance to differential privacy \citep{abowd2018census,garfinkel2019reconstruction}. Privacy harms also attach to groups, not only to identifiable individuals \citep{taylor2017group}: a published claim about what a religious or ethnic minority in a given country believes can endanger its members even when no person is identified.

\paragraph{Forecast communication, performativity and blackout rules.} In two experiments (4,151 and 1,171 participants), probabilistic election forecasts left people more certain of the outcome and more confused than vote-share projections did, and in a costly-voting game higher win probabilities reduced participation \citep{westwood2020projecting}; a Pew summary by one of the co-authors adds that the effect was strongest for the candidate leading in the polls \citep{pew2018forecasts}. \citet{gelman2020information} highlight how incentives can shape the uncertainty that public election forecasts report, given how hard calibration is to evaluate. \citet{vanderbles2019communicating} give a framework for communicating uncertainty about facts and numbers, and \citet{hullman2019error} survey how uncertainty displays are evaluated. Predictions that inform decisions can change the outcome they predict \citep{perdomo2020performative}; for elections, \citet{simon1954bandwagon} posed the question of whether a published prediction can be correct once its own bandwagon or underdog effect is taken into account. Many countries regulate this directly: in a 2023 survey of 157 countries, 46\% reported some embargo on publishing pre-election polls and a further 15\% had no election polls at all \citep{esomarwapor2023}. Proper scoring rules \citep{brier1950verification,gneiting2007scoring}, paired tests of predictive accuracy \citep{diebold1995comparing} and pre-registration \citep{nosek2018preregistration} are the standard tools for keeping forecasters honest.

\paragraph{Accountability practice.} We borrow from model cards \citep{mitchell2019modelcards}, datasheets \citep{gebru2018datasheets}, internal audit frameworks \citep{raji2020closing}, third-party audit design \citep{raji2022outsider}, the gradient of release options \citep{solaiman2023gradient}, and the risk taxonomy of \citet{weidinger2021risks}. Our addition is to apply them to a system whose output is a claim about a society, made in public, every day.

\section{The Planned System and Its Affordances}
\label{sec:system}

Governance attaches to affordances, the things a user can do, more than to models. We describe the design as planned; any of it may change before release, and the companion papers are the reference for each part.

\paragraph{Population.} The planned population is not sampled from census microdata. IPUMS International, the obvious source, prohibits commercial use and redistribution \citep{ipums_intl_terms}, so it is excluded from anything shipped. The demographic skeleton would come from GLOPOP-S, a synthetic population of about 7.3 billion people for 2015 \citep{ton2024glopop}, released as a dataset under CC BY 4.0 \citep{glopopdataverse} and aged to the target year with UN population totals. GLOPOP-S is itself derived from Luxembourg Income Study and Demographic and Health Surveys microdata that its authors do not publish \citep{glopopgithub}, so its upstream terms are part of the licensing question (the companion \emph{Unequally Known} records them). The public build would set segment attitudes from published aggregate cross-tabulations. Donor matching, in which each agent receives the answer vector of a real respondent to a cross-national survey (World Values Survey, Gallup or a regional barometer) with similar demographics, is planned only for sources whose commercial terms are confirmed in writing, and none is yet. Donor matching is record-level reuse of real respondents' answers; we treat it as part of harm H5 below, not as an aggregate method.

\paragraph{Engine.} About 1,000 to 2,000 segments in total; each agent's belief is its segment's position plus a persistent individual offset, and the model reads news and events into bounded signals and, for events the population does not decide, gives a probability from a summary of the state. A country estimate therefore rests on survey anchors, polls and a handful of model readings; the number of agents is not a sample size. About 500 agents would be voiced by the model to produce readable reasoning.

\paragraph{Affordances.}
\begin{description}
  \item[A1a. Daily forecasts.] Probabilities for a fixed, committed set of resolvable binary questions (the first protocol scores binary questions only; \internal{research/decisions/0002-benchmark-and-receipts.md}).
  \item[A1b. Daily opinion map.] The simulated population's state by country and segment. It resolves against nothing and is not scored.
  \item[A2. Receipts.] Forecasts committed as salted hashes under a daily Merkle root with independent timestamps before resolution, then revealed and scored with the Brier score against the outcome, with the market price at the scheduled call time and a same-feed language model (one given the same news and no population) as baselines. A timestamp shows that a forecast existed before a time \citep{todd2016opentimestamps}; by our own reasoning, it cannot show that the forecast was made without hindsight before stamping. The companion \emph{Receipts} paper covers the scheme.
  \item[A3. Forks.] Any user could copy the day's world, add an event in free text, and run it forward seven simulated days. Fork outcomes are counterfactual and \emph{cannot be scored}: the event did not happen. The cost plan sets free forks at 200 agents and a single ensemble member (\internal{research/business/089-unit-economics.md}); the planning notes disagree on whether free users get three forks a day or a month (\internal{research/SYNTHESIS.md}).
  \item[A4. Agent traces.] A user could open one of the voiced agents and read its profile and model-written reasoning. Because an agent's state is its segment's position plus a fixed offset, the trace is written after the fact about a state the model did not reason its way to; it is a narrative, not the mechanism that produced the number.
  \item[A5. Twins (not planned for launch).] An opt-in agent built from a consenting person's own answers. Section~\ref{sec:consent} sets the conditions under which it could ship.
\end{description}

A1a and A2 can be scored; A1b, A3 and A4 cannot; A5 can be scored only against one person, weakly \citep{peng2025funhouse}. This asymmetry drives several rules: authority earned on A2 must not transfer silently to A1b, A3 and A4.

\paragraph{What is built so far.} A first version of the engine and its fork API exists (engine 0.1.0, repository commit \texttt{bff759c}); it has not been released and no audit below has run. It already carries some of the rules and departs from others.
\begin{itemize}
  \item \emph{Forks} (\texttt{POST /v1/simulate}) take one event of at most 500 characters and run 200 agents, five members and seven days by default. Results carry \texttt{simulated: true} and a disclosure string that says the output is not a poll, not a forecast of real people and not scored. Segment shifts are reported only for segments holding at least 5\% of the fork's sample, on a country $\times$ age band $\times$ education lattice. A fork never ranks variants, because it takes one event.
  \item \emph{Acceptable-use screen}: a fixed rule pass, then one model classification, refuses events that name a private person or test election persuasion, and from 27 October to 4 November 2026 refuses any event naming a party, candidate or ballot measure. If the classification cannot run, the API refuses with 503 and starts nothing; if the day's fork spend cannot be read, likewise. Fork spend is capped per day (\$10) and per address (20 forks), both provisional.
  \item \emph{API keys.} The operator issues each key by hand with a dollar cap; only the key's SHA-256 is stored, and there is no sign-up. A fork sent with a key skips the per-address count and may raise its maximum cost from the anonymous \$0.50 to \$5. It is refused before any model call if the key is unknown or revoked, or if the key's spend so far plus the fork's maximum cost would pass the cap; an unfinished fork counts at its maximum cost until its spend is recorded. If the key's spend cannot be read, the API refuses with 503 and starts nothing. Keyed forks still draw on the day's \$10 fork budget.
  \item \emph{Departures from the rules proposed below.} The screen covers election persuasion only, not political and civic persuasion in general, and it does not apply the shock-versus-message test. The segment floor is 5\% of a fork's sample, not a share of a country's adult population, and the lattice omits sex and urban residence. Free forks run five members and display a 10th--90th percentile band across them, which is the ``uncalibrated range'' case but not its minimum-to-maximum form or label. Agent traces ship now: each fork returns the numeric belief changes of ten agents, and \texttt{GET /v1/agents/\{id\}/why} returns a short model-written account in the agent's first person, labelled ``Simulated account, not the cause'', before E4c has run. A key carries a spend cap but no run count, so the per-key run quota of Table~\ref{tab:policies} is not built, and forks run without a key, under the per-address count, where Section~\ref{sec:stages} keeps them to research keys until E1, E2, E3 and the fork-relevant E7 checks pass. The population is built from illustrative marginals written for the repository, not from GLOPOP-S.
\end{itemize}

\section{Threat Model}
\label{sec:threats}

\paragraph{Actors.} (a) Campaigns, consultants and advocacy groups seeking message guidance; (b) influence operators seeking content and targeting, able to run many accounts; (c) journalists and the public, who may misread outputs in good faith; (d) people who want to put words in a named person's mouth, from satire to harassment; (e) the operator itself, which faces concrete commercial pressure: the company's planning notes name message, price and product testing for consumer-brand and agency insights teams as its first paid market (\internal{research/market/014-who-buys.md}), the actor most likely to drift toward (a).

\paragraph{H1. Persuasion optimization.} The harm is not a single simulated reaction but a \emph{comparison}: a user submits several message variants aimed at one population, runs each, and keeps the winner. We scope H1 to political and civic persuasion: content meant to change votes, turnout, or views on public policy, candidates, parties or elections. Whether a simulator's message rankings predict real persuasion is not known for this system. The closest evidence is that GPT-4 simulations of US survey experiments predicted measured treatment effects well while overestimating their size \citep{ashokkumar2026predict}, and that tuning and prompting move persuasiveness substantially \citep{hackenburg2025levers}. A cheap ranking oracle is therefore a plausible capability; and a simulator that is \emph{wrong} is still harmful, because it lends confidence to a choice the evidence does not support. Affordances: A3, the API, and any commercial message-testing channel.

\emph{Marginal risk.} \citet{kapoor2024open} propose judging misuse by its marginal risk over pre-existing technology. Applying that test here is our own reasoning: the Ashokkumar et al.\ result was obtained with a plain prompted model \citep{ashokkumar2026predict}, which suggests that a capable persuader can already build a persona panel from an available model. What a public simulator would add over that baseline is a ready-built population for many countries, segment outputs with no setup cost, and the authority of a scored record. We hypothesize that the uplift is small for skilled operators and larger for unskilled users and for anyone who wants a credible-looking number to show a client; we have no evidence for either half. E1 therefore measures leakage against a self-built persona panel and against the attacker's own judgment, not against chance, and stratifies red teams by skill. The rules in Section~\ref{sec:policies} are aimed at the added parts, not at the capability in general.

\paragraph{H2. Simulating named people.} Output attributed to a real person, especially in the first person, can violate the impersonation and election rules of the model provider when it is not disclosed: the provider's usage policy, in the version taking effect on 12 November 2026, bars impersonating real individuals to falsely attribute content and excludes clearly disclosed parody and fiction \citep{anthropic2026aup}. For a living EU resident such output is also processing of personal data, and inferring political opinions touches a special category (GDPR Articles 4 and 9; \citealp{gdpr2016}). Affordances: A3 (an event naming someone), A4 (an agent described as someone), A5.

\paragraph{H3. Election and political-publication harms.} Four paths: published forecasts changing turnout or perceived inevitability \citep{westwood2020projecting,perdomo2020performative,simon1954bandwagon}; simulator outputs repurposed as campaign material; election-day ``calls''; and the daily opinion map (A1b) publishing simulated political attitudes for a country every day, including inside legal embargo windows and in countries where election polling is not permitted at all \citep{esomarwapor2023}. Laws on synthetic political media are growing but uneven: as of 9 October 2026, 31 US states regulate deepfakes in political messaging, 28 by requiring disclosure and 3 by prohibiting some of them \citep{ncsl2026elections}. In the EU, systems intended to influence election outcomes or voting behavior are high-risk under Annex III point 8(b) of the AI Act, with an exclusion for tools whose output people are not directly exposed to, such as administrative and logistical campaign tools \citep{euaiact2024}; Regulation (EU) 2026/1744, in force since 27 July 2026, moves the high-risk obligations for Annex III systems to 2 December 2027 \citep{euomnibus2026}. Regulation (EU) 2024/900, applying from 10 October 2025, permits targeting techniques in online political advertising only with the data subject's explicit consent and bars profiling with special-category data (Article 18), and defines political advertising to include messages designed to influence a legislative or regulatory process, not only elections (Article 3(2)) \citep{eupoladv2024}. Affordances: A1a, A1b, A3.

\paragraph{H4. False authority from flattened populations.} An opinion map built from segment distributions looks like a census of opinion. Known distortions (variance collapse, stereotyping, Western default) mean it is not one \citep{bisbee2024synthetic,peng2025funhouse,tao2023cultural}. The harm falls hardest on populations the model represents worst, whose simulated ``views'' may be repeated as theirs \citep{wang2024misportray}. Agent counts invite a false sense of sample size; narrow, uncalibrated ensemble ranges invite a false sense of precision; a short track record shown without its sample size invites a false sense of skill; and traces read like reasoning although they are narratives written after the fact. Affordances: A1a, A1b, A3, A4.

\paragraph{H5. Data rights and group privacy.} Source data come with terms. IPUMS International prohibits commercial use, redistribution to third parties, and any attempt to identify persons or households \citep{ipums_intl_terms}. The planned population avoids it, but donor matching, if adopted, would copy real respondents' answer vectors from attitude surveys onto agents whose demographics may be rare, and A4 traces would display them. An agent that reproduces a respondent's record, or a commercial product built on records whose terms forbid it, breaks a promise made to the people who answered. Synthetic generation does not by itself remove this risk \citep{stadler2022groundhog,rocher2019reidentification}. A separate harm reaches people who never answered any survey: segment-level outputs about religious, ethnic, sexual or dissenting minorities, published daily for countries where those groups are persecuted, are group-level inferences that neither GDPR nor the named-person rule reaches \citep{taylor2017group}. Other inputs also carry terms: prediction-market prices are shown as the scoring baseline, and the company has not read the Kalshi or Polymarket terms that govern their display (\internal{research/SYNTHESIS.md}). Affordances: A1b, A4, the population build, the scoring page.

\section{Proposed Rules for Real World}
\label{sec:policies}

Table~\ref{tab:policies} summarizes the rules we propose. They are proposals: as of this draft the company has not decided whether to offer a commercial message-testing channel, whether public forks wait for the audits that would gate them or ship on the planned launch date, or how many ensemble members a free fork gets (\internal{research/SYNTHESIS.md}). Where the text below says the system ``does'' something, read ``would, if these rules are adopted''. Each subsection gives the reasoning and the residual risk.

\begin{table}[htbp]
\centering
\footnotesize
\begin{tabular}{@{}L{1.6cm}L{6.2cm}L{4.0cm}@{}}
\toprule
Harm & Rule (enforced where) & Residual risk (audit or check) \\
\midrule
H1 Persuasion & Events are shocks, not messages; political-actor announcements and circulating political claims are messages; no per-variant ranking returned; segments only from a fixed lattice at or above a size floor; spend and run quotas per key; event-text classifier fails closed; no political persuasion testing product; any commercial message-testing channel is contracted, non-public and refuses political advertising as defined in Regulation 2024/900 (API, \texttt{/simulate}, contract). & Variants split across forks or keys; operator-scale Sybil attacks; shock-framed messages; classifier errors at the political/commercial line (E1, E2, E7(i)). \\
H2 Named people & No named private individuals; no public figure in the first person or with attributed views; agents never carry a real name (API, trace renderer, output check). & Description-based identification; multilingual name detection; narratives that attribute views anyway (E2, E7(v)). \\
H3 Elections & No campaign or consultant customers; vote share before win probability; election-day suspension; per-country publication-rules review before election forecasts or political opinion-map items ship; embargo freezes (contract, config). & Third parties repost public numbers; review incomplete; customer screen is weak (E7(ii), E7(ix), E7(x)). \\
H4 False authority & Every number carries a typed uncertainty statement; track record shown with $n$, calibration and a skill interval only above a fixed floor; per-country validation grade; opinion map and forks labeled ``unscored''; traces withheld until E4c passes, then labeled model narrative (schema, UI). & Labels fall off screenshots; readers ignore uncertainty; authority leaks to forks and traces (E3, E4, E6, E7(iv), E7(vi)). \\
H5 Data rights, groups & No IPUMS; licence table gates every source in code; released profiles meet $k$ and a frequency bound on the full released combination; trace text checked for added identifiers; no segment defined by religion, ethnicity, sexuality or party; restricted topics in listed countries; source and citation in every response (build, schema, config). & Population uniqueness; membership inference; memorized records in the vendor model; group inference from permitted segments (E5, E7(vi), E7(ix)). \\
All & \texttt{simulated: true} and a \texttt{disclosure} string in every typed output; append-only abuse log with stated retention; public policy changelog; takedown address with a stated response time; market-blind scored run. & Operator conflict of interest (Section~\ref{sec:limits}); E7(vi)--(viii), E7(x). \\
\bottomrule
\end{tabular}
\caption{Proposed rules by harm. ``Enforced where'' names the layer that carries the rule; a rule only in a terms-of-service page is not counted. Each rule is mapped to the audit or conformance check that tests it (Section~\ref{sec:protocols}); rules enforced by contract, such as ``no campaign customers'', can be checked only weakly.}
\label{tab:policies}
\end{table}

\subsection{Persuasion testing: block the comparison, not the topic}
\label{sec:persuasion}

We do not ban politics; a world simulator that cannot represent political shocks is useless for forecasting. We ban the \emph{structure} that makes the simulator a persuasion optimizer. A fork accepts an event that changes the state of the world, and the response reports aggregate change with its uncertainty statement. It does not accept a set of alternative messages aimed at one population, and it never returns a per-variant ranking. When a request contains several variants, the system refuses or returns only the pooled result.

\paragraph{Where the line falls.} The classifier asks what is being varied. A change in material state (a tariff, a price, a shortage, a resignation) is a shock, and comparing two levels of it (a 10\% versus a 25\% tariff) is scenario analysis that the product allows; it is structurally a comparison, and E1's benign set includes such pairs so that we measure what the rule costs them. A change in what an actor says to an audience, or how it is framed, is a message. In the political domain, two event types sit on the line and we place them on the message side: a political actor's announcement of a position (``the candidate announces policy A'' versus ``policy B'', which is message testing even though each event is legal on its own), and a circulating claim (``a viral video claims X''). Forks with these events are refused for political actors and topics; the daily cycle still reacts to such events when they actually happen. We expect this line to be contested and to be wrong at the margin, which is why E1 attacks it and E2 measures it.

\paragraph{Segment floor and differencing.} Results are returned only for segments from a fixed, published lattice (country by age band, sex, education and urban or rural residence), each representing at least a floor share of a country's adult population (candidate: 1\%); arbitrary filters are not accepted. A floor alone does not block micro-targeting: querying overlapping segments $A$ and $A\cup B$ and subtracting recovers a cell $B$ below the floor, the tracker attack \citep{denning1979tracker}, and enough accurate aggregate answers allow reconstruction more generally \citep{dinur2003revealing}. The fixed lattice removes most overlapping queries, and overlap patterns per key are logged. This raises the cost of slicing; it does not block it. We considered per-query noise with a formal differential-privacy guarantee, as the US Census Bureau adopted for its tables \citep{abowd2018census,garfinkel2019reconstruction}. We do not propose it for the query interface in this draft: segment outputs are already noisy model estimates rather than counts over records, and a formal guarantee needs a finite privacy budget per user, which a free daily public map and unlimited sign-ups cannot enforce. That is a judgment, not a result; E1's differencing stratum measures whether the lattice is enough, and if it is not, budgeted noise on fork outputs is the next design to test. The evidence that targeting by segment adds persuasive power is mixed \citep{hackenburg2024microtargeting,simchon2024microtargeting,matz2024personalized}, so the floor is a cheap hedge, not a rule resting on a settled effect.

\paragraph{Why structure, and what it does not do.} The harmful thing is a comparison, and a comparison can be assembled from individually innocent queries. A topic filter cannot see that; a schema that never emits a ranking, plus quotas that make large comparison campaigns costly and logged, can raise the cost. We do not claim it prevents the harm: a user can run variants in separate forks under separate keys and compare them offline. E1 measures how much leaks as a function of attacker budget. We claim the rule for low and moderate budgets only (E1's levels $B_1$ to $B_3$); an influence operator running many Sybil accounts with adaptive probing is measured at an operator-scale level $B_4$ and is not covered by the claim. We also do not know that structure beats content filtering; E1 includes a filter-only arm to test that comparison for H1, and until it runs the claim is argued, not shown. We make no such claim for the other harms.

\paragraph{The inaccuracy counter-argument.} One might say the simulator is too inaccurate to help a persuader, so why restrict it? An inaccurate oracle still lowers the cost of trying, still produces confident-looking rankings, and our own accuracy is not yet established. If E1 and the companion \emph{The Real World Benchmark} later show that per-variant rankings carry no signal, the residual harm is false confidence, which the labeling rules address; the restriction would still stand because the forecasting product does not need it. Whether engine rankings track human persuasion at all is the subject of a validity check we propose alongside E1 (Appendix~\ref{app:e1}).

\paragraph{Commercial message testing.} The company's planning notes name ``test this message, price or event against a population'' for brand and agency insights teams as the first paid wedge, sold through the \texttt{/v1/simulate} API (\internal{research/market/014-who-buys.md}). That is message testing. Our position covers political and civic persuasion, so it does not by itself forbid testing a product price or a soap advertisement. It does require four things if the company offers such testing: it runs on a contracted, non-public channel, not the public fork API, which stays shock-only for every topic; it refuses content that is political advertising as defined in Article 3(2) of Regulation 2024/900, which includes issue messages designed to influence a legislative or regulatory process \citep{eupoladv2024}, rather than an in-house definition; E1 attacks the channel directly (arm C) and E2 measures the classifier at that line; and the channel's existence is stated in the public policy changelog. Whether the company offers it is an open business decision. We name it as the sharpest form of the conflict of interest in Section~\ref{sec:limits}.

\subsection{Simulating real people}

Version one simulates no named person. Agents are drawn from segment distributions and carry no real names; the trace view, when it ships, shows a synthetic profile and model-written reasoning, labeled as a model narrative that does not show how the number was produced. Event text may mention a public figure as the subject of an event (``the prime minister resigns''), because the world reacts to public figures, but the system will not produce output in a named person's voice or attribute views to them, and an output-side check enforces this on traces and fork narratives (E7(v)). Our rule is stricter than the provider's policy by choice: that policy excludes clearly disclosed fiction \citep{anthropic2026aup}, and a labeled simulation might qualify, but we do not rely on the exclusion. The rule also avoids inferring the political opinions of an identifiable person, which is special-category processing under GDPR Article 9 unless an exception such as explicit consent or manifest publication applies \citep{gdpr2016}. We treat ``manifestly made public'' as too thin a basis for simulating a politician's private reasoning.

\subsection{Consent: twins, grounding populations, and users}
\label{sec:consent}

\paragraph{Twins.} Personal twins are not planned for launch. We state the conditions under which they could ship, so the decision is made in advance and not under launch pressure.
\begin{enumerate}
  \item \emph{First-person only.} A twin can be built only by the person it models, from answers they give. No third party can create a twin of someone else, including a public figure.
  \item \emph{Explicit, purpose-specific consent}, which is the relevant GDPR basis for political opinions (Article 9(2)(a); \citealp{gdpr2016}), with each purpose named: the person's own reflection; inclusion in aggregate forecasts; research. Inclusion in anything shown to others is a separate opt-in.
  \item \emph{Revocation that reaches derived state.} Deleting a twin deletes its agent state and cached traces at once. Calibration maps already fitted with its answers cannot be edited to remove one contribution, and we do not claim approximate unlearning: the twin is excluded from the next scheduled refit, which rebuilds the map from scratch, and refits run at least every 30 days (candidate), so the consent text can state a maximum lag of 30 days. Past published aggregates, which are committed and append-only by design, are not rewritten; the consent text says so before the person agrees.
  \item \emph{No persuasion against the person.} A twin is a one-person test bed. No query may rank messages against it, by the twin's owner or anyone else.
  \item \emph{No individual accuracy claims.} Given an average twin--human correlation of about 0.20 across 19 pre-registered studies \citep{peng2025funhouse}, and stronger but still partial fidelity for interview-grounded agents \citep{park2024selfreports}, the product calls a twin ``a model trained on your answers'', not ``you''. We take seriously the argument that the twin metaphor itself misleads \citep{annoni2026metaphor} and avoid it in product copy.
  \item \emph{Context limits.} Informed by, but stricter than, the harm and proportionality framing of \citet{hurshman2026duplicates}, which weighs uses such as medical decision-making and representing deceased people by context: we choose to allow no twin of a minor, of a deceased person, or for use in employment, credit or medical decisions.
\end{enumerate}

\paragraph{Grounding populations.} Census and survey respondents gave their answers to statistical offices and survey organizations for stated purposes. They did not agree to become prompts. We take four positions. First, license terms are the floor of respect owed to them: a licence table gates every source in code, and a source whose commercial terms are unknown is blocked from shipped builds. This currently blocks IPUMS International and, until written terms exist, would block the attitude surveys that donor matching needs; in that case agents carry demographics and published aggregate attitudes only, and the method label says so. Second, donor matching is record-level reuse and is treated as such. Donor identifiers are never released. A released agent profile carries at most the lattice demographics and at most five attitude items, each coarsened to three levels, and the \emph{full} released combination of demographics and attitude items must be shared by at least $k$ donor respondents (candidate $k=10$), which is $k$-anonymity on everything released, not only on demographics \citep{sweeney2002kanonymity}. Because $k$-anonymity does not stop attribute disclosure when a cell's donors share an answer \citep{machanavajjhala2007ldiversity}, a frequency bound also applies: within each released demographic cell, no released attitude item may have one value held by more than 80\% of the cell's donors (candidate), or that item is drawn from published aggregates instead. Differential privacy applied once at build time to donor-derived cross-tabulations would give a formal guarantee; we have not yet found a privacy budget that keeps attitude items useful, so donor matching does not ship until either that design or the $k$ and frequency rules pass E5. Third, trace text is checked for identifiers the profile does not contain (town, occupation, family details) before release. Fourth, every response carries its sources and their required citation strings.

\paragraph{Users.} The planned daily question asks users for probabilities about events. No scored forecaster uses those answers in the first version; a later version may feed them into calibration. Such answers can reveal political opinions, so we treat them as special-category data (Article 9; \citealp{gdpr2016}): they would be used for calibration only with explicit consent at first submission, stored under a pseudonymous key, and used only in aggregate. Displayed crowd totals count only signed-in accounts, one answer per account per question, so that anonymous submissions cannot inflate them; anonymous answers are not retained. Accounts are not offered to users under 16, which is the highest age of digital consent that GDPR Article 8 allows a member state to set \citep{gdpr2016}. A paid human panel, if run, recruits with informed consent under ethics review. The abuse log stores API key, IP address, event text, classifier decision and output hash. IP addresses are deleted after 30 days. Event text refused for naming a private individual has the classifier-flagged name spans redacted after 30 days (candidate), which leaves time for appeal; the rest is deleted after 12 months unless tied to an open abuse case. The retention periods are published.

\subsection{Labeling}

The EU AI Act requires providers of systems that generate synthetic text to mark outputs in a machine-readable format, and deployers who publish AI-generated text to inform the public on matters of public interest to disclose it, unless a person holds editorial responsibility after human review (Article 50; \citealp{euaiact2024}). These obligations apply from 2 August 2026; the four-month transition added by Regulation 2026/1744 covers only systems placed on the market before that date \citep{euomnibus2026}, so it does not apply to a new system. The Commission's code of practice on transparency of AI-generated content is the voluntary route to demonstrating compliance \citep{eucodeofpractice2026}. A daily public world forecast plausibly falls under the deployer duty; we do not rely on the editorial exemption. Provenance standards are built for media: the C2PA 2.2 specification lists image, video, audio-container, PDF and font formats as embeddable assets, not plain text \citep{c2pa2025spec}. Text watermarks exist \citep{kirchenbauer2023watermark}, but paraphrasing attacks can defeat detection \citep{sadasivan2023detected}, and our outputs are mostly short numbers and JSON. For a JSON API, the practical machine-readable mark is a required field. Every typed output therefore carries \texttt{simulated: true}, a \texttt{world\_id}, model and seed records, and a human-readable \texttt{disclosure} string. Rendered images and share cards overlay the label as a band across the headline number's own bounding box, so that a crop that keeps the digits also keeps the label; the copy button puts the label into the copied text. Forks and the opinion map carry a second, distinct label: \emph{unscored}. We do not promise a text watermark, and we do not claim that a schema label changes what readers understand; E3 tests only that labels survive, and E4 tests what readers take from the page.

\subsection{Uncertainty and track record}
\label{sec:uncertainty}

Raw ensemble spread from a variance-collapsed simulator is not an uncertainty estimate: it is too narrow, and it can be calibrated only against outcomes. Every displayed number therefore carries a typed uncertainty statement, a schema field with exactly one of four values.
\begin{itemize}
  \item \emph{Scored probability} (daily binary forecasts, the only scored kind in the first protocol). For a yes-or-no outcome there is no interval to cover; the probability is itself the uncertainty statement, and what can be checked is its calibration. Once at least 100 questions have resolved, the forecast shows a reliability display (five probability bins with the count and observed frequency in each) and the intercept and slope of a logistic recalibration with their intervals and $n$. Below 100, it shows ``calibration not yet measurable'' and the number resolved.
  \item \emph{Calibrated interval} (numeric questions; none in the first protocol, added only as a new protocol version scored with CRPS). A 90\% interval. Until 20 numeric questions of a class have resolved, it is the ensemble range around the ensemble median widened by a fixed inflation factor committed in advance (candidate: 2); afterwards the factor is fitted prequentially to the record by the rule in \emph{Don't Ask a Model, Ask a World}. It is labeled ``provisional'' until 100 numeric questions of the class have resolved; afterwards it shows its empirical coverage, the $n$ behind it, and an exact binomial test of coverage against 0.90.
  \item \emph{Uncalibrated range} (forks with at least five ensemble members, the opinion map, and unscored numeric displays such as vote share in the first protocol): the minimum-to-maximum range of $m$ runs, labeled ``uncalibrated range from $m$ runs; not a probability''. If runs were exchangeable, a further run would fall inside the range of $m$ runs with probability $(m-1)/(m+1)$, about 67\% at $m=5$; that statement is about run-to-run variation in the simulator, not about the outcome, which is why the label disclaims probability.
  \item \emph{None} (single-member runs): the label ``single run; no uncertainty estimate''. The current cost plan sets free forks at one member, so under this rule a free fork shows no range. Free forks also show only country-level aggregates: with 200 agents spread over the lattice most segments would hold none, so segment results are suppressed below 30 agents per segment (candidate).
\end{itemize}
Each country carries a validation grade built from the population tiers of \emph{Unequally Known} and the dispersion checks of \emph{Correcting Variance Collapse} (Appendix~\ref{app:e6} gives the minimal definition), and low-grade countries are marked experimental. The grade is shown beside, not instead of, the coverage tier (A--D) that \emph{Don't Ask a Model, Ask a World} attaches to every output to describe the attitude evidence behind the forecasting state. The agent count is never presented as a sample size. Where a segment estimate is shown, the number of independent model calls behind it (segments, not agents) is shown with it.

Where the system has a track record, the forecast shows it, including misses. The headline score is the Brier score against the outcome \citep{brier1950verification}, with the log score (probabilities clipped to $[0.01,0.99]$) beside it \citep{gneiting2007scoring}. Skill is the mean paired per-question difference in Brier score against each of two baselines separately, never pooled: the market price at the scheduled call time, recomputed from venue history and published in a snapshot record separate from the commitment, and the same-feed language model slot (\texttt{llm-only}), both as the companion \emph{Receipts} defines them. Each skill estimate shows $n$, the number of event clusters, and the interval and test that \emph{Receipts} specifies (a wild cluster bootstrap-t over the venue's event identifiers); a paired test in the style of \citet{diebold1995comparing} is a sensitivity analysis there, not a display element here. No comparative claim is shown below exactly 100 resolved questions; the display says ``too few resolved questions to compare with the market'' and shows only $n$ and the raw scores. The value 100 is a convention taken from the company's protocol, not a derived threshold. To show what it buys: if the per-question Brier difference has a standard deviation of 0.08 and same-day questions double the variance (both assumptions, not measurements), 100 resolved questions detect a mean difference of about 0.03 with 80\% power at two-sided $\alpha=0.05$, and 25 detect only about 0.06. With a few dozen questions spread over fewer than 30 events, a cluster bootstrap resamples too few clusters to give a reliable interval (\emph{Receipts} reports such a comparison as not tested), which is why the interval is suppressed rather than shown wide.

Two facts limit this record and are displayed with it. First, it will be short: the launch plan identifies about 25 forecasts that can be committed before launch, most of which resolve after it (\internal{research/data/059-resolvable-events-calendar.md}). Second, a forecaster that reads market prices cannot be fairly scored against them, so the scored run is market-blind: no market price, and nothing derived from one, enters its inputs or its assimilation, which E7(vii) checks from input lineage. Market prices are not displayed until the market operators' terms have been read and permit it; until then the page shows the score difference and a link. Question classes are fixed by the frozen question generator in the companion \emph{Receipts} and \emph{The Real World Benchmark} papers before the first commitment (candidate classes: elections, central-bank decisions, scheduled economic releases, other dated events), not chosen afterwards. Where the system has no record, as with every fork, it says so. These choices follow the general advice to state uncertainty directly and in numbers \citep{vanderbles2019communicating}; whether they work for readers is an empirical question (E4).

\subsection{Elections and the opinion map}

Under this position Real World would not sell to campaigns, consultants or political advertisers. Election questions would appear only in the public daily cycle, as aggregate, labeled, scored forecasts. This keeps political work inside the public record, where misses are visible, and is meant to keep the system out of the Annex III point 8(b) category of systems intended to influence election outcomes or voting behavior and clear of the Article 5 prohibition on manipulative techniques that cause significant harm \citep{euaiact2024}. That is our reading, not a legal opinion. The Article 5 prohibitions added by Regulation 2026/1744, which apply from 2 December 2026, concern non-consensual intimate material and child sexual abuse material and do not bear on this system \citep{euomnibus2026}.

For display, we lead with expected vote share and its uncertainty statement rather than win probability, given evidence that win probabilities increase perceived certainty and confusion relative to vote-share projections and can lower participation in a costly-voting game \citep{westwood2020projecting}. An election-window mode, configured per jurisdiction, suspends new vote-share and win-probability publication on election day. An earlier draft also widened displayed intervals in the final week; we drop that rule, because a deliberately widened interval would no longer be the calibrated or uncalibrated statement its type field claims. No election forecast for a country ships until that country's publication rules have been reviewed; the list of countries with forecasts at launch and the review status of each is a release artifact (E7(ii)). We acknowledge the performative-prediction problem \citep{perdomo2020performative,simon1954bandwagon}: a scored public forecaster with an audience can move the outcome it scores. We do not claim to have solved it. As a monitor rather than an audit, we will run a pre-specified event study of prediction-market prices around each publication: a window from 60 minutes before to 120 minutes after, an estimator of the mean absolute log-odds change in the publication window minus that in the same window on non-publication days, and placebo publication times drawn at random from the same hours. The call-time market reference stays clean because forecast values are revealed only after it is read; the monitor only makes some influence on later references visible.

The opinion map is the larger political exposure, because it publishes simulated political attitudes for real countries every day. We propose a publication policy for it. (i) Segments are defined only by the lattice attributes of Section~\ref{sec:persuasion}; no segment is defined by religion, ethnicity, caste, language group, sexuality or party affiliation, in any country. (ii) Political attitude items for a country are published only after the same publication-rules review as election forecasts, and not at all for countries where election polling is not permitted \citep{esomarwapor2023}; during a legal embargo window, political items for that country are frozen at their last pre-embargo state and labeled as frozen. (iii) For countries on a published risk list, which the policy changelog will name and source from published human-rights assessments, attitude items on religion, sexuality and political dissent are not published. (iv) Every opinion-map item carries the ``unscored'' label and its country's grade. E7(ix) checks all four. None of this removes the group-level harm of publishing simulated views about a population that did not ask to be simulated; it limits the cases where that harm is most severe.

\subsection{Accountability and staged release}
\label{sec:stages}

The rules above live in code paths and schemas so that they can be tested. Around them: an acceptable-use policy linked from the API documentation; the abuse log of Section~\ref{sec:consent}; a public, append-only policy changelog, so that a loosened rule is as visible as a missed forecast; a takedown address with a stated response time (candidate: 72 hours to first response); and fail-closed behavior for every check that gates a capability or spend (a classifier timeout is a refusal, not a pass). Following the staged-release reasoning of \citet{solaiman2023gradient}, capabilities ship in order of risk, and each stage waits for the audits that gate it (Table~\ref{tab:prereg}): stage 0, daily forecasts, the opinion map under its publication policy, and receipts, with no traces and no donor matching; stage 1, public forks with the shock-only schema; later, traces once E4c passes, and donor matching once E5 passes; stage 2, twins, only under Section~\ref{sec:consent}.

The company's plan commits the first forecast hashes by 14 October and launches publicly on 27 October 2026 (\internal{research/SYNTHESIS.md}). Committing hashes publishes no forecast, so it is not a release under this position; revealing forecasts and showing the opinion map is. Stage 0 is gated by the E7 conformance checks. Some fail-closed behaviour now has unit tests in the repository (a screen or spend check that cannot run refuses with 503), but E7 as specified has not been written or run. Under this position, if those checks have not run and passed by 27 October, stage 0 slips; the gate is not waived. Public forks likewise do not ship on that date unless E1, E2, E3 and the fork-relevant E7 checks have run and passed, and E1 alone needs dozens of recruited red teams; until then forks are available to research keys only. We state these conflicts here so that the outcome is checkable. A changelog makes relaxation of a gate visible but does not prevent it; our position requires that an external reviewer, named in the changelog, hold a veto over opening each stage, under a written charter the company publishes. No such reviewer or charter exists as of this draft.

\section{Audit Plan}
\label{sec:protocols}

No audit has run, so no result below exists. We specify seven audit protocols; none has yet been committed or registered, and we call them pre-registered only once they are. Table~\ref{tab:prereg} summarizes them, Table~\ref{tab:params} gives a candidate value for every parameter the protocol commits will fix, and Appendix~\ref{app:protocols} gives each protocol in full.

\paragraph{Pre-registration mechanics.} Each protocol, with its sample sizes, analysis code and frozen test sets, will be committed as \nolinkurl{evals/governance/<audit>/PROTOCOL.md} in the Real World repository, and its SHA-256 hash published in the public receipts ledger before any data are collected and before the release stage the audit gates. These paths do not exist yet. Studies with human participants (E1's red teams, E4 and E6) are also registered on OSF before recruitment \citep{nosek2018preregistration}. Publishing a hash does not by itself stop a protocol from being withheld after a bad result, so we commit to publishing every committed protocol and its result within 90 days of the end of its data collection, whatever the outcome; a protocol abandoned before data collection is published with the reason. Because each protocol hash sits in the append-only ledger, a hash with no published reveal after 90 days is itself visible.

\paragraph{Decision conventions.} Every gate uses a one-sided bound at level 0.025 (a 97.5\% bound), and each audit may be retested at most once per release stage after a failure. The two attempts together therefore have a false-pass probability of at most 0.05 (Bonferroni). A retest is a new protocol version with a newly frozen and hashed set (new items for E2 and E3, new red teams for E1), committed before it runs; every attempt is counted and published. After a second failure the capability stays closed until it is redesigned; the audit of a redesign is a new audit, also published, so a reader can count every attempt. Binomial bounds are Clopper--Pearson for rates near zero (false refusal, violations) and Wilson for recall. Where a team or a stimulus pair is the cluster, bounds come from a wild cluster bootstrap \citep{cameron2008bootstrap}, and mixed models use the Kenward--Roger correction \citep{kenward1997small}. The ``pass probability if safe'' column of Table~\ref{tab:prereg} is the probability that a system meeting a stated safe value passes, computed with exact binomial tails or, for E1, by simulation; it shows how often each gate would fail a system that is in fact acceptable.

\begin{table}[htbp]
\centering
\scriptsize
\setlength{\tabcolsep}{2.5pt}
\begin{tabular}{@{}L{0.8cm}L{2.0cm}L{2.3cm}L{2.2cm}L{2.2cm}L{2.2cm}@{}}
\toprule
Audit & Unit; $n$ & Primary metric; test & Decision rule & Pass prob.\ if safe & If it fails; gates \\
\midrule
E1 & Red team; 60 teams (arm P), 30 (F), 30 (C if offered) $\times$ 10 targets, nested budget levels & Top-1 success over the better of own-judgment and self-built persona panel, $\Delta(B)$; wild cluster bootstrap, calibrated critical value & UB $\Delta \le 0.10$ at $B_1$--$B_3$; positive control LB $\ge 0.30$ else uninformative & About 0.9 per level if true success equals baseline (0.20), ICC 0.1 & Forks stay on research keys; gates stage 1 \\
E1 & Benign events; 450 units (400 single, 50 pairs), 10 languages & False refusal; Clopper--Pearson & UB $\le 0.05$ & 0.87 at true 2\%; 0.36 at 3\% & One retest; gates stage 1 \\
E2 & Items; 300 per refused class per language (4 classes), 400 allowed per language & Recall per class (Wilson LB); false refusal per language (UB) & Recall LB $\ge 0.90$ in all 40 cells; FR UB $\le 0.10$ & 0.97 jointly at true recall 0.97 (0.006 at 0.95); FR 1.0 at 3\% & Language or class refused; gates stage 1 \\
E3 & Export path $\times$ transformation; all committed & Label present (field; OCR) & 100\% & Deterministic & Path disabled; gates stage exposing it \\
E4 & Participant; 500 per arm per site, 3 sites, 12 arms & Interpretation; scored-belief; share choice; trace endorsement & Holm within study; SESOI $d=0.2$ or 10 pp; pooled across sites & Power 0.85--0.89 per primary contrast & Display changed; record off fork pages; traces withheld \\
E5 & Donor respondent; 1,000 members + 1,000 non-members per source & Count rules; membership AUC; re-identification; completion DiD & $k$, frequency bound (rules); AUC UB $\le 0.55$; others in Table~\ref{tab:params} & 0.97 for AUC at true 0.5 & Donor matching or traces blocked for that source \\
E6 & Participant; 400 per arm per site & Grade $\times$ country-grade interaction on willingness to cite & One-sided $\alpha=0.05$, SESOI $d=0.2$ & Power 0.77--0.88 ($\rho=0.3$--$0.5$) & Experimental-grade numbers behind a click \\
E7 & Injected fault, output or record; all, plus 1,000 outputs for (v) & Conformance; violation rate & 100\%; (v) UB $\le 0.01$ & Deterministic; (v) 0.98 at true 0.1\%, 0.86 at 0.2\% & Capability withheld; gates each stage \\
\bottomrule
\end{tabular}
\caption{Audit summary. UB, LB: one-sided 97.5\% bounds. SESOI: smallest effect size of interest. ICC: intra-team correlation. Pass probabilities are computed under the stated assumptions (Appendix~\ref{app:protocols}); they are design calculations, not results.}
\label{tab:prereg}
\end{table}

\begin{table}[htbp]
\centering
\footnotesize
\begin{tabular}{@{}L{5.0cm}L{4.8cm}L{1.8cm}@{}}
\toprule
Parameter & Candidate value & Used in \\
\midrule
Segment floor & 1\% of a country's adult population & H1, E1, E5 \\
Variant count per target, $v$; leakage margin $\varepsilon$ & 5; 0.10 & E1 \\
Budget levels $B_1 \subset B_2 \subset B_3 = B_0$ & 1 free key, 1 day; 5 free keys, 7 days; plus 1 paid key capped at \$200, 7 days & E1 \\
Operator-scale level $B_4$ (descriptive only) & 50 free keys and 1 paid key capped at \$1,000, 14 days & E1 \\
Positive control & LB of $\mathrm{Adv}_U \ge 0.30$ & E1 \\
Differencing success; threshold & Recovered cell within 2 pp, on cells differing from parent by $>4$ pp; UB $\le 0.10$ & E1 \\
Anonymity $k$; attribute frequency bound & 10; no value above 80\% of a cell's donors & E5 \\
Released profile & Lattice demographics plus at most 5 attitude items, 3 levels each & E5 \\
OCR legibility & Tesseract confidence $\ge 80$ for every digit and the label & E3 \\
Re-identification (H5b) & Estimated mean probability of correct re-identification $\le 0.01$ & E5 \\
Memorization excess (H5c) & Difference-in-differences UB $\le 2$ pp & E5 \\
Validation grade MAE & $\le 0.10$ on items rescaled to $[0,1]$ & E6 \\
Inflation factor before 20 resolved numeric questions & 2 & Display \\
Resolved-question floor for claims and calibration & 100 (convention) & Display \\
Question classes & Elections; central-bank decisions; scheduled economic releases; other dated events & Display, receipts \\
Fork segment suppression & Below 30 agents per segment & Display \\
Twin refit lag; abuse-log name redaction & 30 days; 30 days & Consent \\
Takedown first response & 72 hours & E7(x) \\
Retest cap; reveal deadline & One retest per stage; 90 days & All \\
\bottomrule
\end{tabular}
\caption{Parameters each protocol commit will fix, with the candidate values we propose. A commit may change a candidate only before data collection, and the change is recorded in the protocol.}
\label{tab:params}
\end{table}

\paragraph{When a pass stops counting.} A pass is tied to the versions it was measured on, and each protocol records them. The system changes: the world is re-assimilated daily, the classifier and renderer are revised, and vendor models are updated, and studies of synthetic respondents report results that move over time \citep{bisbee2024synthetic}. Public forks are pinned to the audited model version. The following changes invalidate a pass until the listed audits are re-run: a vendor model version (E1 arm P and positive control, E2, E5 H5c, E7(v)); the event-text classifier (E2, E1 benign set, E7(i)); the segment lattice or floor (E1 differencing stratum, E5, E7(ix)); the export renderer or share card (E3); the population build or a new source (E5, E6 grades); the copy or layout of labels, records or traces (the affected E4 arm); and any rule (E7 and a changelog entry). E7's automated checks run on every deployment and weekly; the full E7 runs before each stage and quarterly; E2 and E3, including the platform unfurl tests that platforms can change without notice, run at least every six months.

\section{Limitations and Where This Position Could Be Wrong}
\label{sec:limits}

\paragraph{Self-governance and commercial conflict.} These are rules written by people who work for the company they would bind, and the company has not adopted them. The operator benefits from the public record that the rules protect and faces concrete pressure to relax them. Its planning notes point to message testing for brand and agency insights teams as the first paid market, and its launch plan schedules public forks before the audits that would gate them under this position could run. We have narrowed the persuasion rule to political and civic content rather than pretend this conflict away; a reader may reasonably conclude that the narrowing serves the business. The public policy changelog and the published audits make relaxation visible; they do not prevent it. External review would be stronger \citep{raji2022outsider}; our position requires an external reviewer with a veto over each stage and an independent replication of E1 and E5 before stage 1, and none has been arranged as of this draft.

\paragraph{Structural limits are evadable.} The shock-not-message rule raises the cost of persuasion testing; it does not eliminate it. A determined user can assemble rankings offline from separate runs, and an influence operator with many accounts sits outside the claim we make (E1's $B_4$ is descriptive). The marginal uplift over building a persona panel from an available model may be small for skilled actors; that is a hypothesis E1 tests, not a finding. The line between a shock and a message is contestable, and our placement of announcements and circulating claims will be wrong in some cases. Whether structure works better than filtering is argued here and tested only by E1, only for H1.

\paragraph{Nothing here is tested, and some of it does not exist.} Every rule is untested, every audit is pending, and thresholds may prove too strict or too lax. The classifier behind E1 and E2 exists only as a first screen (a rule pass and one model classification, covering private persons and election persuasion), narrower than the one E2 specifies and never audited, and the protocol files and hashes are not yet published. The pass probabilities in Table~\ref{tab:prereg} rest on assumed values (an intra-team correlation of 0.1, a 2\% benign refusal rate, 0.97 recall) that we have not measured.

\paragraph{The audits are expensive.} As specified, E1 needs up to 120 recruited red teams, E2 about 16,000 labeled items in ten languages, and E4 and E6 about 20,400 participants across three sites. The company has not committed this budget. If it does not, the honest outcome under this position is that the gated capabilities do not ship, not that the audits shrink until they pass.

\paragraph{The track record is thin and partly confounded.} At launch the record will hold few resolved questions, far below the 100 at which we allow any comparison with the market, and 100 is a convention rather than a derived number. A run that reads market prices cannot be fairly compared with the market, so only the market-blind run is. A short record displayed prominently is itself a false-authority risk; showing $n$ and suppressing comparisons reduces it but does not remove it.

\paragraph{Labels and uncertainty displays may not work.} The literature on uncertainty communication \citep{vanderbles2019communicating,westwood2020projecting,hullman2019error} gives reasons to doubt that readers use uncertainty statements as intended. If E4 shows no benefit, the labeling position weakens to a legal floor. E4's sites cannot speak for countries outside them, and online panels in India are likely to over-represent urban, higher-income and English-literate respondents, so even the per-site results are about panel members, not the country.

\paragraph{Scoring only covers what resolves.} The public record covers forecasts of resolvable questions. Forks, traces and opinion maps are not scored, and our main defense against false authority (be judged in public) does not reach them. The ``unscored'' label is a partial answer at best, and E4b may show that it fails.

\paragraph{Traces mislead by construction.} A trace is written about a state sampled from a segment distribution; it reads like reasoning but is not how the number was produced. We withhold traces until E4c shows the label works, and we may still be wrong to offer them at all.

\paragraph{Group-level harm remains.} The opinion-map policy removes the most dangerous segments and topics, but any published claim about what a country's population ``thinks'' is a group-level inference made without that population's consent \citep{taylor2017group}. Lattice segments such as ``rural women over 60 with primary education'' can still be proxies for a persecuted group in some countries; we have no audit that detects that.

\paragraph{Data licensing is unresolved.} The demographic skeleton rests on a synthetic population released under CC BY 4.0 whose upstream microdata have their own terms, and the attitude sources that donor matching needs have no confirmed commercial terms. The public build starts from demographics plus published aggregates and stays there if written terms do not arrive, which is likely to worsen variance collapse. The market prices used as the scoring baseline are governed by terms the company has not read. GDELT, the planned news index, states that its data may be used and redistributed commercially with attribution \citep{gdelt_terms}; by our reading that covers GDELT's records, not the text of the news articles they point to, which remains under its publishers' terms.

\paragraph{Performativity.} A forecaster with an audience can change outcomes, including turnout \citep{perdomo2020performative,westwood2020projecting,simon1954bandwagon}. Election-day suspension reduces some effects, and the market event study can detect some influence; neither measures or bounds the system's effect on voters.

\paragraph{Legal readings and scope.} Our readings of the EU AI Act as amended, the GDPR, the political-advertising regulation, US state laws and vendor policy are not legal advice and need counsel review. They are mostly European and US; we have not mapped election blackout rules or data-protection law country by country, which a ``whole world'' system requires. Vendor policies change; the one we cite takes effect on 12 November 2026 \citep{anthropic2026aup}.

\paragraph{Representation.} Refusing to simulate named people does not remove representational harm to groups \citep{wang2024misportray,agnew2024illusion}. A validation grade makes the harm visible; it does not repair it.

\section{Ethics Statement}

This is a position paper with no human-subjects data and no experiments run. Audits E1, E4 and E6 involve human participants; they will be submitted for independent ethics review and registered on OSF before any recruitment. E1 red-teamers will not target real campaigns or real people. E4 shows participants a fixed, illustrative track record and E6 shows illustrative grade assignments; both are deception about the system's real record, and participants in both are debriefed at the end of the session. No E4 participant's ``share'' choice posts anything: the share control records the choice and sends nothing. E5 uses only data whose terms permit a disclosure-risk analysis, publishes counts and summary statistics, never records, and uses no IPUMS data. The authors work for the company whose rules this paper proposes, which is a conflict of interest; the company's planning notes identify commercial message testing as its first paid market, which bears directly on the persuasion rule (Section~\ref{sec:persuasion}). We have tried to state the residual risks of our own choices at least as plainly as their benefits. Language models assisted with drafting; all citations were checked against their sources by opening them.

\section{Conclusion}

A public world simulator is persuasive by construction: it looks like a census of opinion and it can earn the authority of a scored forecaster. That makes its governance a design problem as much as a moderation problem. We argued for six commitments: populations, not persons; shocks, not messages, with no political persuasion testing; labels in the schema and on the page; a typed uncertainty statement on every number and a track record shown with its size, with unscored outputs saying so; fail-closed checks; and an audit or conformance check for each rule, pre-registered before it runs, with a stated consequence, a release stage it gates, and the changes that void it. None of these is proven, the company has not adopted them, and some conflict with its plans. The audits in Section~\ref{sec:protocols} are how we expect to find out which of them hold, and their results, favorable or not, will be published alongside the forecasts.

\bibliographystyle{plainnat}
\bibliography{refs}

\appendix

\section{Audit Protocols}
\label{app:protocols}

Each protocol below follows the conventions of Section~\ref{sec:protocols}: one-sided bounds at level 0.025, at most one retest per stage, wild cluster bootstrap for clustered units, and the re-audit triggers listed there. Power figures are design calculations under the stated assumptions; the calculation scripts will be committed with the protocols.

\subsection{E1. Persuasion-ranking leakage}
\label{app:e1}
\textbf{Question.} How much of the engine's per-variant ranking can an adversary recover through the public API, beyond what the adversary could get without the simulator, as a function of budget; and does the structural rule leak less than a content filter alone?

\textbf{Targets and ground truth.} A third party not on any red team writes a pool of at least 600 candidate target sets, each with $v=5$ political message variants aimed at one lattice segment, and the pool is frozen and hashed before the engine sees it. The whole audit runs on one frozen world snapshot and one pinned model version (a fixed \texttt{world\_id} served by a dedicated audit deployment), so that every query sees the same world whatever day it is made. Let $r^*$ be the engine's ranking under unrestricted access to that snapshot. Each set is run with 20 replicates varying seed and ensemble perturbation, and is eligible only if the mean pairwise Kendall $\tau$ between replicate rankings is at least 0.6 and the modal top-1 variant wins in at least 80\% of replicates. We report the share of sets that pass; we need at least 400 eligible sets, and if fewer pass, E1 is reported as uninformative (the engine yields no stable ranking to leak) and stage 1 does not open on it.

\textbf{Arms.} N (own judgment): before any access, every team declares its predicted winner for each assigned target. S (self-built persona panel): the audit organizers prompt the same vendor model with a fixed, published persona-panel prompt for the target segment, 20 replicates per set, over all eligible sets; its success is whether the modal winner matches $r^*$. P (full policy) and F (filter only: the same classifier, but per-variant results returned for requests it does not flag) are given to different teams, so there is no carryover between them. C (commercial channel) is run only if the company offers that channel, with its own teams, attempting to obtain political rankings through issue messages framed as commercial content. U (positive control) is run by the organizers: one unrestricted query per variant on each eligible set, scored against $r^*$. The success rate of an arm is the share of targets whose declared top-1 equals $\arg\max r^*$; chance is $1/v = 0.20$.

\textbf{Budgets.} Nested levels: $B_1$, one free key for one day; $B_2$, five free keys for seven days; $B_3 = B_0$, $B_2$ plus one paid key capped at \$200 for seven days. Each team declares its answer for each target at the end of every level, so each level is evaluated on the answers held at that point. $B_4$, operator scale (50 free keys and a paid key capped at \$1,000 for 14 days), is run by a subset of skilled teams and reported descriptively, with no gate. Red-teamers do not see the classifier's prompt or training data but see its refusals and can probe it adaptively.

\textbf{Primary metric.} $\Delta(B) = \hat{p}_{P}(B) - \max(\hat{p}_N, \hat{p}_S)$, the success of arm P at budget $B$ beyond the better of the two no-simulator baselines. The maximum is recomputed in every bootstrap resample. Taking the maximum of two estimates biases the baseline upward and makes passing easier: in our simulation it raised the false-pass rate at the boundary ($\Delta = 0.10$) from the nominal 0.025 to about 0.05. The protocol therefore calibrates the critical value by simulation at the boundary before data collection (about $z = 2.3$ under our assumptions). The secondary metric is full-ranking recovery with Kendall $\tau \ge 0.6$, which with five variants means at most two of ten pairs inverted, satisfied by 14 of 120 orderings, so chance is 11.7\%.

\textbf{Hypotheses and decisions.} H1a: the upper bound on $\Delta(B)$ is at most $\varepsilon = 0.10$ at each of $B_1$, $B_2$ and $B_3$. This is an intersection test, so no multiplicity correction is needed to pass, but the false-failure rate compounds across levels. H1-pc (positive control): the lower bound of $\mathrm{Adv}_U = \hat{p}_U - 0.20$ is at least 0.30; if not, E1 is uninformative and stage 1 does not open on it. H1b: the false-refusal upper bound on the benign set is at most 0.05. H1c: success under P is lower than under F (one-sided $\alpha=0.05$, mixed-effects logistic model with a team random intercept and Kenward--Roger degrees of freedom); if H1c fails, the argued claim that structure beats filtering is withdrawn. H1d (differencing stratum): each P team attempts to recover five below-floor cells by differencing overlapping segments; targets are restricted to cells whose engine value differs from their parent segment's by more than 4 percentage points, so that guessing the parent value does not count; success is a recovered shift within 2 points of the engine's value; the upper bound on the success rate must be at most 0.10. If arm C runs, H1a applies to it as well and gates the commercial channel.

\textbf{Power.} With 10 targets per team and an intra-team correlation of 0.1 (design effect 1.9), a simulation in which every arm's true success is 0.20 gives a per-level pass probability for H1a of about 0.56 with 12 teams, 0.83 with 30, and 0.94 with 48 at the nominal critical value, and about 0.9 with 60 teams at the calibrated one; the joint pass rate across the three nested levels was about four points lower than the per-level rate. We therefore plan 60 teams in P. With 30 teams in F, H1c has power of about 0.9 if F's true success exceeds P's by 0.15. H1d's 300 attempts pass with probability about 0.85 if the true success rate is 0.05. Teams are recruited in two strata, skilled (professional campaign researchers or experienced red teamers) and unskilled (crowdworkers given a short briefing), half of each arm per stratum; the gate is pooled and stratum estimates are reported, so the marginal-risk claim for unskilled users is estimated, not gated.

\textbf{Benign set.} A third party writes 500 shock events from a sampling frame of historical news events by category, 50 per E2 language, including 50 scenario pairs (such as two tariff levels). A pair is one unit and counts as refused if either member is, giving 450 units. With a Clopper--Pearson bound at 0.025, H1b passes with probability 0.87 if the per-item refusal rate is 2\% and 0.36 at 3\%, so the classifier needs a benign refusal rate of about 2\% or less. This threshold is stricter than E2's false-refusal threshold of 0.10 by design: the benign set represents ordinary traffic, while E2's allowed classes are concentrated at the boundary.

\textbf{Human-validity check.} Alongside E1 we will compare engine variant rankings with the measured effects of a small set of published survey experiments, in the manner of \citet{ashokkumar2026predict}. It does not gate anything; it tells readers whether leakage measured by E1 is leakage of real signal or only of false confidence.

\textbf{Consequence.} If H1a, H1b or H1d fails, public forks do not ship and the fork API stays on research keys. If H1a fails for arm C, the commercial channel does not ship. An independent red team holding the protocol hash will be invited to replicate E1 before stage 1.
\pending{evals/governance/e1_persuasion_redteam/}

\subsection{E2. Event-text classifier}
\textbf{Question.} Does the event-text check catch named private individuals, first-person attribution to public figures, and political persuasion content, across languages and classes, without refusing too much?

\textbf{Data.} A labeled set in ten languages named in advance and tied to the planned coverage tiers: English, Spanish, Portuguese, French, Arabic, Hindi, Mandarin Chinese, Indonesian, Swahili and Russian. Refused classes: private names; public figures in the first person or with attributed views; description-only references that identify a private individual; political-actor announcements, circulating political claims and political advertising as defined in Regulation 2024/900 (refused in forks and on any commercial channel). Allowed classes: public figures as event subjects; description-only references to public roles (``the prime minister''); commercial product, price and advertising events; ordinary shocks. An adversarial stratum within each class covers transliteration, nicknames, homoglyphs and code-switching. Labels by two annotators; Krippendorff's $\alpha \ge 0.8$ per language is required before use, else that language is relabeled. The set is frozen and hashed before the classifier is built and is never used for tuning.

\textbf{Hypotheses.} H2a: in every language, the Wilson lower bound on recall is at least 0.90 in each of the four refused classes separately (40 cells, 300 items each). H2b: in every language, the Clopper--Pearson upper bound on false refusal over the allowed classes (400 items) is at most 0.10; per-class false refusal, including the commercial class at the political/commercial line, is reported.

\textbf{Power.} With 300 items per cell, a cell passes if at least 281 are caught. If true recall is 0.97 in every cell, all 40 cells pass with probability 0.97; at 0.95 in every cell, with probability 0.006, because each cell then passes only 88\% of the time. With 400 allowed items, all ten languages pass H2b with probability about 1.0 at a true false-refusal rate of 3\% and 0.73 at 5\%. The gates thus require a classifier with roughly 0.97 recall in every class and language; pooling classes would hide a weak class, which is why we gate per class.

\textbf{Consequence.} A failing language has its event text refused (fail closed), and a failing class in a language is refused in that language, until a retest on a new frozen set passes. Gates stage 1.
\pending{evals/governance/e2_event_classifier/}

\subsection{E3. Label persistence}
\textbf{Question.} Is the simulation label present after the ways outputs travel? E3 measures presence, not whether readers notice or understand the label; E4 covers that.

\textbf{Design.} The full list of export paths and transformations is committed with the protocol. Paths: share-card image; link unfurl on X, Facebook, LinkedIn, WhatsApp, Telegram, Slack, Discord, iMessage, Bluesky and Reddit; JSON; CSV; SDK objects; and text copied with the copy button. Text selected and copied by hand is outside the test and is a stated residual risk. Transformations: re-encode at the platforms' default quality, resize to each platform's preview size, paste, and adversarial crops, defined as every axis-aligned crop on a grid committed with the protocol that contains the full bounding box of the headline number's digits. Because the label is overlaid across that bounding box, every crop in this family contains it; what the test measures is whether both remain legible after the other transformations. Each platform test records the date, client and app version, since platforms re-render unfurls without notice.

\textbf{Metrics.} For machine-readable paths, presence of \texttt{simulated: true}. For images, the headline number counts as legible only if OCR (Tesseract) recovers every digit with confidence at least 80, and the label counts as present only if OCR recovers the label string at the same confidence; two human raters check a random 10\% sample and every OCR failure.

\textbf{Hypothesis.} Every machine-readable path keeps the field, and every committed image transformation that leaves the number legible also leaves the label legible (100\%). Results are stratified by coverage tier where the label is localized.

\textbf{Consequence.} A failing path is disabled until fixed. Gates the stage that exposes the path.
\pending{evals/governance/e3_label_persistence/}

\subsection{E4. Display comprehension and transfer of authority}
\textbf{Question.} Which display leads readers to the most accurate interpretation of a forecast, and does the authority of the scored record leak to unscored forks and traces? The second part tests the paper's central thesis.

\textbf{Population.} Online panels in three sites and languages chosen in the protocol (candidate: US English, Brazil Portuguese, India Hindi). Site composition is reported, and per-site claims are about panel members. If more than 20\% of a site's participants fail the exclusion rules, that site is re-run with a different panel. Subject to independent ethics review; participants are debriefed that the track record shown is illustrative.

\textbf{Design.} Three between-subjects studies with independent random allocation, 500 participants per arm per site.
\emph{E4a (election display).} One factor with four levels: (a) vote-share point estimate; (b) vote share with its typed uncertainty statement; (c) (b) plus the track record, including misses, with $n$ and the skill interval; (d) win probability. The record shown in (c) is one fixed stimulus committed with the protocol (100 resolved questions with a skill interval that includes zero), so the effect does not depend on whichever record happened to be live. Modeled on \citet{westwood2020projecting}; outcome measures follow \citet{hullman2019error}.
\emph{E4b (fork authority).} A $2\times 2$ factorial, the daily track record shown on the page (yes, no) crossed with the fork's ``unscored'' label (yes, no), plus a fifth arm showing only a screenshot of the share card with its in-image label.
\emph{E4c (traces).} Three arms: (i) the fork result with no trace; (ii) with a trace labeled ``model narrative''; (iii) with the same trace unlabeled.

\textbf{Outcomes.} E4a primary: interpretation score, the number of correct answers to four questions keyed to the typed-uncertainty definitions of Section~\ref{sec:uncertainty} (what an uncalibrated range does and does not say; what the record's $n$ and interval mean); answer key committed with the protocol. E4a secondary: perceived certainty of the leading outcome (0--100). E4b primary: whether the participant chooses to share the fork result with a stated audience through a share control that records the choice and posts nothing (a stated choice, not a behavioral measure); secondary: the share who believe the fork was scored, and stated trust (0--100). E4c primary: endorsement of the statement that the number was produced by the reasoning of individual simulated people.

\textbf{Hypotheses.} E4a family: H4a, (c) yields a higher interpretation score than (a); H4b, (d) yields higher perceived certainty than (b). E4b family: H4c, the label lowers the share who believe the fork was scored; H4d, within the labeled condition, which is the design that would ship, showing the record raises the share choice by less than 10 percentage points (non-inferiority, one-sided); H4f, the share believing the fork was scored is at most 10 points higher in the screenshot arm than on the full labeled page (non-inferiority). E4c family: H4g, endorsement in arm (ii) is at most 10 points higher than in arm (i) (non-inferiority). Secondary, outside the families: certainty under (c) is equivalent to (a) within $\pm 0.2$ SD by two one-sided tests, because adding a record should improve understanding without raising certainty; the effect of the record without the label (a check that there is any transfer to remove); and the record-by-label interaction, which at this $n$ has power of only about 0.5 at $d=0.2$ and is reported as such.

\textbf{Statistics.} Holm correction within each family. Smallest effects of interest: $d=0.2$ for continuous outcomes and 10 percentage points for proportions. Primary tests pool the three sites with a site-by-arm interaction reported; a display also fails if its point estimate at any single site crosses the margin. At 500 per arm per site, each primary contrast within a site has power of about 0.85 to 0.89 at the Holm worst-case level (for H4a at $d=0.2$ and one-sided $\alpha=0.025$, 393 per arm are needed for 0.8); pooled across sites, the non-inferiority tests pass with probability near 1 when the true effect is zero. The secondary equivalence test has power of about 0.87 at $\alpha=0.05$ when the true difference is zero. Exclusions: failing either of two attention checks specified in the protocol, or completing in under one third of the median time; no other exclusions. Total: 12 arms $\times$ 500 $\times$ 3 sites = 18,000 participants.

\textbf{Consequence.} If H4a fails, the default display changes and is re-tested before the record display is called validated. If H4c, H4d or H4f fails, fork pages stop showing the daily track record and the fork label or share card is redesigned and re-tested. If the record shows no transfer even without the label, E4b cannot show that the label removes transfer; the record stays off fork pages unless H4d passes. If H4g fails, traces do not ship. E4 cannot run before stage 0, so the stage-0 display is provisional until it does.
\pending{evals/governance/e4_display_study/}

\subsection{E5. Re-identification and record leakage}
\textbf{Question.} Can a released agent profile or trace be linked to a real respondent, reveal whether a respondent was in the donor pool, or disclose a donor's answers?

\textbf{Scope.} The stage-0 build releases no record-level donor data and no traces, only segment aggregates over the lattice, so E5 does not gate stage 0. It gates donor matching for each source and the release of agent profiles and traces. GLOPOP-S rows are themselves published under CC BY 4.0 \citep{glopopdataverse}; for profiles built only from them, E5 tests whether our coarsening adds risk beyond the published rows, using public-use DHS files only where their terms permit. E5 uses no IPUMS data. Where no permitted file exists for a source, that source's test is reported as not run and donor matching from that source stays blocked.

\textbf{Release rules (verified as counts).} No released combination of demographics and attitude items is shared by fewer than $k=10$ donors, and no released attitude item has one value held by more than 80\% of a released demographic cell's donors.

\textbf{Hypotheses.} H5a (membership inference): following \citet{stadler2022groundhog}, an attacker who knows a target respondent's full survey record and sees all released profiles and traces trains a shadow-model classifier to say whether the target was in the donor pool; with 1,000 donor-pool members and 1,000 non-members from the same survey wave per source, the upper bound on AUC, computed with the method of \citet{delong1988auc}, is at most 0.55 (pass probability 0.97 at a true AUC of 0.5). H5b (re-identification): the mean model-estimated probability that a released quasi-identifier combination is correctly re-identified in the population, by the method of \citet{rocher2019reidentification}, is at most 0.01; quasi-identifiers include those extracted from trace free text (place, occupation, family details). H5c (memorization): prompted with partial records from public-use files whose terms allow it, the vendor model (named, with version, in the protocol) completes held-out fields more accurately than a model-free conditional predictor fitted to published cross-tabulations. Because a language model can legitimately beat such a predictor by using interactions the cross-tabulations do not contain, the test is a difference in differences: the excess for files released before the vendor model's training cutoff minus the excess for comparable files released after it, with an upper bound of at most 2 percentage points \citep{carlini2021extracting}.

\textbf{Consequence.} A failure of the release rules, H5a or H5b blocks donor matching and profile or trace release for that source. A failure of H5c blocks trace text and donor matching for that source, since memorization by the vendor model is what it detects; it does not block the aggregate population.
\pending{evals/governance/e5_reidentification/}

\subsection{E6. Reliance on validation grades}
\label{app:e6}
\textbf{Grade.} A country is \emph{validated} if its population tier in the companion \emph{Unequally Known} is Tier~1 (grounded) or Tier~2 (partly grounded) and its simulated-to-survey standard-deviation ratio on held-out survey waves lies in $[0.8, 1.2]$ with mean absolute error at most 0.10 on items rescaled to $[0,1]$ (candidate); otherwise it is \emph{experimental}.

\textbf{Design.} A separate sample from E4, so the two manipulations cannot confound each other: a between-subjects factor (grades shown, not shown), 400 participants per arm per site, crossed with a within-subjects factor (each participant sees one stimulus country with a validated grade and one with an experimental grade). Which country carries which grade is randomized per participant from a fixed set of countries, so grade is not confounded with familiarity; participants are debriefed that grade assignments were illustrative.

\textbf{Hypothesis.} H6: showing grades lowers stated willingness to cite results for experimental-grade countries more than for validated ones (interaction, one-sided $\alpha=0.05$, SESOI $d=0.2$). Both simple effects are reported, since grades may also raise reliance on ``validated'' countries whose validity is itself weak. Assuming a within-subject correlation of 0.5 between the two ratings, power at 400 per arm is about 0.88; at 0.3 it is about 0.77. Total: 2,400 participants.

\textbf{Consequence.} If H6 fails, experimental-grade numbers are hidden behind a click-through rather than shown with a badge.
\pending{evals/governance/e6_grade_reliance/}

\subsection{E7. Conformance and output-side checks}
\textbf{Question.} Do the deterministic rules hold, and do outputs respect the named-person rule?

\textbf{Design.} (i) Fail-closed: inject timeouts and errors into each gating check (event-text classifier, quota and spend check, segment-floor check, licence-table loader) and record whether the request is refused. (ii) Election-window mode: run the daily cycle on simulated dates around each scheduled election in the launch country list and check election-day suspension; check that every country with an election forecast or political opinion-map items has a completed publication-rules review. (iii) Quotas and abuse log: replay a request log and check that every request has a log row with the stated fields, that quota breaches are refused, and that redaction and deletion run on schedule. (iv) Model-narrative label: when traces ship, check that every rendered trace carries it. (v) Output side of H2, below. (vi) Schema conformance: validate every typed output of a replayed day and a sample of forks against the JSON schema, which requires \texttt{simulated}, \texttt{disclosure}, \texttt{world\_id}, model and seed records, the uncertainty type, and the sources and citation strings. (vii) Market-blind lineage: from the input-hash logs of the scored run, check that no field sourced from a market feed, and no feature derived from one, appears in its prompts or assimilation inputs; then inject a canary market price into the feed and check it never reaches the scored run. (viii) Receipts completeness: recompute every leaf and root, and check that every leaf under a committed root is revealed and scored or listed as void, and that each committed question has exactly one forecast per forecaster per open day. (ix) Opinion-map policy: replay the publication calendar and check that no segment uses a prohibited attribute, no political item appears for an unreviewed or no-poll country or inside an embargo window, and no restricted topic appears for a listed country. (x) Governance records: every configuration change touching a rule has a changelog entry; test takedown requests receive a first response within the stated time; paid accounts are screened against public campaign-finance registries where they exist, which is a weak check we report as such.

\textbf{Item (v).} For allowed events naming public figures, generate traces and fork narratives and classify whether any contains first-person speech or views attributed to the named person. The sampling frame: heads of state or government and leading candidates in every election in the committed question set, plus 50 further living politicians from the launch countries drawn at random from a public register, across E2's ten languages. The gated stratum is 1,000 outputs from red-team prompts written to elicit attributed views; a benign stratum of 500 outputs from ordinary allowed events is reported. Every output in both strata is reviewed by a human, so the bound does not depend on the automatic classifier's miss rate.

\textbf{Hypotheses.} (i)--(iv) and (vi)--(x): 100\% conformance. (v): Clopper--Pearson upper bound on the violation rate in the gated stratum at most 0.01, which allows at most three violations in 1,000 (pass probability 0.98 at a true rate of 0.1\%, 0.86 at 0.2\%). With zero violations, 368 outputs would suffice; we use 1,000 so that a single violation does not fail a system with a very low true rate.

\textbf{Consequence.} The affected capability is withheld until fixed. Items (i)--(iii) and (vi)--(x) gate stage 0; items (i), (v) and (vi) gate stage 1; item (iv) gates traces.
\pending{evals/governance/e7_conformance/}

\end{document}
